
Let @@M0@@ be the homomorphism so defined. Denoting by @@M1@@ the ideal @@M2@@ of @@M3@@, one has @@M4@@. Keeping the notation of Lemma 2.1.4, one obtains a homomorphism of abelian groups
@@M5@@
that associates to the class in @@M8@@ of every invertible @@M6@@-module @@M7@@ the element @@M9@@ of @@M10@@.

Lemmas 2.1.2 and 2.1.4 immediately give:

\sgaStableTarget{sga6.e10.prop.2-2-1}\textbf{Proposition 2.2.1.} \uline{For every noetherian scheme} @@M11@@ \uline{and every} @@M12@@ (@@M13@@\uline{), the composite map}
@@M14@@
\uline{defined by}
@@M15@@
\uline{is a polynomial map of degree} @@M16@@.

\sgaStableTarget{sga6.e10.sec.2.3}\subsection*{2.3.}
If @@M17@@ is a proper scheme over a field @@M18@@ and @@M19@@ is the structural morphism, then the homomorphism
@@M20@@
associates to the class @@M21@@ of every coherent @@M22@@-module @@M23@@ the Euler--Poincaré characteristic @@M24@@ of @@M25@@ on @@M26@@. From 2.2.1 one immediately concludes:

